\begin{tikzpicture}[
  font=\scriptsize,
  kepala/.style={align=center, font=\tiny\bfseries},
  label/.style={anchor=west, align=left, font=\scriptsize},
  isi/.style={align=center, font=\tiny\bfseries}
]
  % Ukuran: kolom label 3,3 cm, delapan kolom masing-masing 1,08 cm,
  % baris judul 0,9 cm, sembilan baris isi masing-masing 0,66 cm.
  \fill[praditagreen!15] (0,0) rectangle (11.94,-0.9);

  % Sel berisi, diwarnai sebelum garis kisi digambar.
  \fill[green!10]  (3.30,-0.90) rectangle (4.38,-1.56);
  \fill[blue!8]    (4.38,-1.56) rectangle (5.46,-2.22);
  \fill[blue!8]    (4.38,-2.22) rectangle (5.46,-2.88);
  \fill[blue!8]    (5.46,-2.88) rectangle (6.54,-3.54);
  \fill[blue!8]    (6.54,-3.54) rectangle (7.62,-4.20);
  \fill[blue!8]    (7.62,-4.20) rectangle (8.70,-4.86);
  \fill[gray!35]   (10.86,-4.86) rectangle (11.94,-5.52);
  \fill[red!22]    (10.86,-5.52) rectangle (11.94,-6.18);
  \fill[orange!12] (8.70,-6.18) rectangle (9.78,-6.84);
  \fill[orange!12] (9.78,-6.18) rectangle (10.86,-6.84);

  % Garis kisi.
  \foreach \y in {0,-0.9,-1.56,-2.22,-2.88,-3.54,-4.20,-4.86,-5.52,-6.18,-6.84}
    \draw[gray!60] (0,\y) -- (11.94,\y);
  \foreach \x in {0,3.30,4.38,5.46,6.54,7.62,8.70,9.78,10.86,11.94}
    \draw[gray!60] (\x,0) -- (\x,-6.84);
  \draw[thick] (10.86,0) -- (10.86,-6.84);
  \draw[thick] (0,-6.18) -- (11.94,-6.18);

  % Judul kolom: elemen target, ditambah kolom Dihapus.
  \node[kepala] at (1.65,-0.45) {Baris: \textit{baseline}\\Kolom: target};
  \node[kepala] at (3.84,-0.45) {ERP\\keuangan};
  \node[kepala] at (4.92,-0.45) {SIMRS\\terpadu};
  \node[kepala] at (6.00,-0.45) {Platform\\integrasi};
  \node[kepala] at (7.08,-0.45) {Platform\\data};
  \node[kepala] at (8.16,-0.45) {\textit{Landing}\\\textit{zone}};
  \node[kepala] at (9.24,-0.45) {MPI};
  \node[kepala] at (10.32,-0.45) {Akses\\terpusat};
  \node[kepala] at (11.40,-0.45) {Dihapus};

  % Judul baris: elemen baseline, ditambah baris Baru.
  \node[label] at (0.08,-1.23) {ERP keuangan};
  \node[label] at (0.08,-1.89) {Lima SIMRS lokal};
  \node[label] at (0.08,-2.55) {SIM Klinik};
  \node[label] at (0.08,-3.21) {Integrasi \textit{point-to-point}};
  \node[label] at (0.08,-3.87) {23 laporan \textit{spreadsheet}};
  \node[label] at (0.08,-4.53) {Server lokal rumah sakit};
  \node[label] at (0.08,-5.19) {Antrean Bandung};
  \node[label] at (0.08,-5.85) {Surveilans infeksi Bogor};
  \node[label, font=\scriptsize\bfseries] at (0.08,-6.51) {Baru};

  % Isi sel.
  \node[isi] at (3.84,-1.23) {Tetap};
  \node[isi] at (4.92,-1.89) {Diganti};
  \node[isi] at (4.92,-2.55) {Digabung};
  \node[isi] at (6.00,-3.21) {Diganti};
  \node[isi] at (7.08,-3.87) {Diganti};
  \node[isi] at (8.16,-4.53) {Diganti};
  \node[isi] at (11.40,-5.19) {Sengaja};
  \node[isi] at (11.40,-5.85) {Tidak\\[-2pt]sengaja};
  \node[isi] at (9.24,-6.51) {Baru};
  \node[isi] at (10.32,-6.51) {Baru};

  % Keterangan.
  \node[font=\tiny, text width=11.8cm, align=center] at (5.97,-7.27)
    {Tetap: ada di kedua sisi. Diganti atau digabung: elemen lama punya padanan
     baru. Baru: dibangun, dibeli, atau berlangganan. Dihapus: diperiksa, sengaja
     atau tidak.};
\end{tikzpicture}
